\documentclass[11pt]{article}

% Standalone mirror of ../THEORY.md. amsmath/amssymb only; no external figures,
% no external packages, so it compiles with a bare LaTeX installation.
% (Compile-untested in the build environment; written to compile with pdflatex.)

\usepackage{amsmath}
\usepackage{amssymb}

\newcommand{\code}[1]{\texttt{#1}}
\DeclareMathOperator{\STFT}{STFT}
\DeclareMathOperator{\iSTFT}{iSTFT}
\DeclareMathOperator{\Rea}{Re}
\DeclareMathOperator{\clip}{clip}
\DeclareMathOperator{\Var}{Var}
\DeclareMathOperator{\SE}{SE}
\newcommand{\eps}{\varepsilon}
\newcommand{\vt}{\tilde v}
\newcommand{\at}{\tilde a}
\newcommand{\norm}[1]{\lVert #1 \rVert}
\newcommand{\inner}[2]{\langle #1,#2\rangle}
\newcommand{\keptbar}{\overline{\langle a\rangle}_{\text{kept}}}
\newcommand{\dropbar}{\overline{\langle a\rangle}_{\text{dropped}}}

\title{Theory Notes --- Direction 06:\\
Robust Training under Stem Bleed}
\author{SingNet project}
\date{Frozen 2026-07-13}

\begin{document}
\maketitle

\begin{abstract}
This note derives the mathematics implemented in \code{singnet/\{data,losses,analysis,
train\}/} for Direction 06: how a controlled stem-bleed corruption of the training
targets hurts a compact mask-based separator, and why a cheap classical defense (trimmed
loss) can recover some of the loss. The crux is the pair (\S3,\,\S5): \S3 gives a
falsifiable closed-form null --- the clean SI-SDR of a model that perfectly learns the
corrupted target, $\STFT$-exact and machine-checked; \S5 derives the mechanism by which
trimming, run \emph{outside} its proven sample-level regime, selects the effectively
cleanest chunks under \emph{uniform} corruption. Every \S3 number is machine-checked
against \code{singnet/analysis/bleed.py} on synthetic stems.
\end{abstract}

Notation. Clean vocal $v$, accompaniment $a$, mixture $x=v+a$; complex STFTs $V,A,X=V+A$.
The separator predicts a soft mask $M\in[0,1]^{F\times T}$ and estimates
$\widehat V=M\odot|X|\,e^{i\angle X}$ (masked mixture magnitude, \emph{mixture phase});
$\hat v=\iSTFT(\widehat V)$. Bleed level $\eps\in\{0,0.05,0.15,0.30\}$.

\section{Corruption as continuous-target label noise}

Applied at the stem level, at load time, before augmentation, to training targets only:
\begin{equation}
\boxed{\ \vt = v + \eps\,a,\qquad \at = (1-\eps)\,a\ }
\end{equation}
a \emph{redistribution} of a fraction $\eps$ of the accompaniment into the vocal target.

\paragraph{Mixture invariance.} $\vt+\at = v+\eps a+(1-\eps)a = v+a = x$: the corrupted
stems still sum to the true mixture, so the network input is untouched and the model learns
$x\mapsto v+\eps a$ --- to leave $\eps a$ in its vocal estimate. This is the defining
property of real bleed and of \code{SDXDB23\_Bleeding}. The map is deterministic (no RNG),
so wiring it in cannot perturb any augmentation stream --- the corrupted arm changes only
the target's content (unit-tested).

\paragraph{Loudness (pre-registered).} $\norm{\vt}^2=\norm{v}^2+2\eps\inner{v}{a}+
\eps^2\norm{a}^2$ grows with $\eps$; we do \emph{not} renormalize (real bleed adds energy;
the input standardization is mixture-side; random-gain already randomizes levels;
renormalizing would confound $\eps$ with attenuation). The prediction line (\S3) uses the
same un-normalized construction, so the comparison is internally consistent.

\section{The corrupted-optimal L1-magnitude model}

With loss $\mathcal L=\operatorname{mean}\big|M\odot|X|-|\vt|\big|$ and corrupted target
magnitude $|\vt|=|V+\eps A|$, the L1 objective is separable across TF bins. At a bin, the
minimizer of $|M|X|-t|$ over $M\in[0,1]$ with $t=|V+\eps A|$ is
\begin{equation}
\boxed{\ M^\star=\clip\!\Big(\tfrac{|V+\eps A|}{|V+A|},0,1\Big)\ }
\quad\Longrightarrow\quad M^\star|X|=\min\big(|V+\eps A|,|V+A|\big).
\end{equation}
Since $\eps<1$ the corrupted target carries less accompaniment than the mixture, so
$|V+\eps A|\le|V+A|$ on almost all bins (clip inactive) and the optimum reproduces the
\emph{corrupted-target} magnitude. Expanding,
\begin{equation}
|V+\eps A| = |V| + \eps\,\underbrace{\Rea\!\big(A\,e^{-i\angle V}\big)}_{\Delta}
    + O(\eps^2),\qquad |\Delta|\le|A|,
\end{equation}
so the L1-optimal masked magnitude is $|V|+\eps\Delta+O(\eps^2)$: the model is trained to
output the clean vocal magnitude \emph{plus} a systematic $\eps$-scaled accompaniment term.
It leaks $\eps a$ by construction; the residual against $|V|$ is $\eps\Delta\ne0$.

\paragraph{Phase caveat.} The mask wears the mixture phase $\angle(V+A)$, not
$\angle(V+\eps A)$, so even the magnitude-perfect estimate is not the complex $V+\eps A$.
The idealized $\hat v=v+\eps a$ (\S3) is thus a magnitude-\emph{and}-phase ideal upper-bounding
the mask model; the mixture-phase penalty is $\eps$-independent (already present at
$\eps=0$), so the \emph{degradation curve} $s(0)-s(\eps)$ cancels it and isolates the bleed
cost. The dose--response curve, not the absolute SI-SDR, is the clean observable.

\section{The prediction line $P(\eps)$}

The idealized model outputs $\hat s=v+\eps a$; its clean SI-SDR against $s=v$ is closed-form.
With $\rho=\inner{a}{v}/(\norm{a}\norm{v})$,
\begin{equation}
\alpha=\frac{\inner{\hat s}{s}}{\norm{s}^2}
   =1+\eps\,\frac{\inner{a}{v}}{\norm{v}^2}
   =1+\eps\,\rho\,\frac{\norm{a}}{\norm{v}},
\end{equation}
and the noise is
\begin{equation}
\alpha v-\hat s=(\alpha-1)v-\eps a
   =-\eps\Big(a-\tfrac{\inner{a}{v}}{\norm{v}^2}v\Big)=-\eps\,a_\perp,
\end{equation}
so \emph{only the component of $a$ orthogonal to $v$ survives as SI-SDR noise} --- the
$v$-parallel bleed is absorbed by the scale-invariant $\alpha$. With
$\norm{a_\perp}^2=\norm{a}^2(1-\rho^2)$,
\begin{equation}
\boxed{\ P_{\text{track}}(\eps)=10\log_{10}
   \frac{\alpha^2\,\norm{v}^2}{\eps^2\,\norm{a}^2\,(1-\rho^2)}\ }
\qquad(\text{exact projection form}).
\end{equation}

\paragraph{Orthogonality and the $-20\log_{10}\eps$ slope.} When $v\perp a$ ($\rho=0$),
$\alpha=1$ and
\begin{equation}
P_{\text{track}}(\eps)=10\log_{10}\frac{\norm{v}^2}{\eps^2\norm{a}^2}
   =\underbrace{-20\log_{10}\eps}_{\text{universal slope}}
   +\underbrace{10\log_{10}\tfrac{\norm{v}^2}{\norm{a}^2}}_{\text{per-track anchor}},
\end{equation}
a $-20\log_{10}\eps$ line ($\approx6.0$ dB per halving of $\eps$) anchored by each track's
energy ratio. \code{prediction\_line} returns both the exact and the orthogonal columns;
their difference measures departure from orthogonality. Machine-checked: orthogonal
$v=[1,1,1,1],a=[1,-1,1,-1]$ gives $P(0.1)=20.0$ dB with exact$=$orthogonal; non-orthogonal
$v=[1,0],a=[1,1],\eps=0.5$ gives $\alpha=1.5,\rho^2=\tfrac12$, so exact
$=10\log_{10}9=9.54$ dB vs orthogonal $=10\log_{10}2=3.01$ dB.

\paragraph{Reading the gap.} Let $s(\eps)$ be the measured clean curve. $s\approx P$:
faithful learning of the corrupted conditional (no implicit denoising). $s>P$: implicit
robustness (the model beats the ``leak exactly $\eps a$'' estimator). $s<P$: corruption
also destabilizes optimization. If the curve is flat, $P(\eps)$ must sit well below it,
turning ``flat'' into a quantified robustness claim. $P$ is computed from clean stems only.

\section{Trimmed estimators, honestly compared}

ITLM (Shen \& Sanghavi, 2019) alternates selecting the $\alpha=1-q$ fraction of lowest-loss
samples and retraining on them; it \emph{provably} recovers ground truth in GLMs when a
fraction of samples are arbitrarily corrupted and \emph{the rest are clean} --- a
\emph{sample-level} premise. Arpit (2017): DNNs learn general patterns first and memorize
noise later, so early high loss $\approx$ likely-corrupted.

Under Direction 06's \emph{uniform} bleed every target is corrupted by the same $\eps$:
there are \emph{no} clean samples, so ITLM's premise fails and a null is pre-registered.
What trimming can still do is select the \emph{less-corrupted} chunks --- those whose
$\eps a$ happens to be quiet. Whether that has teeth depends on whether per-chunk loss
correlates with per-chunk corruption magnitude; \S5 shows it does, through accompaniment
energy. We run the classic trick outside its proven regime and report the mechanism.

The wrapper keeps $\lceil(1-q)B\rceil$ chunks per step (MASTER\_PLAN \S3.2; for $q=0.30,
B=16$ this is $\lceil11.2\rceil=12$); $q=0$ reduces to the base loss.

\section{Selection mechanism under uniform bleed}

Fix a converged \emph{general} model $f$ (clean vocal structure learned, chunk noise not
memorized). Its per-chunk L1-mag loss on corrupted chunk $i$ is, by the \S2 expansion,
\begin{equation}
\ell_i(f)=\operatorname*{mean}_{TF}\big|f(|X_i|)-|V_i+\eps A_i|\big|
   \approx\eps\operatorname*{mean}_{TF}|\Delta_i|,
\qquad \Delta_i=\Rea\!\big(A_i e^{-i\angle V_i}\big).
\end{equation}
By Cauchy--Schwarz $\operatorname{mean}|\Delta_i|$ is bounded by, and tracks, the chunk's
accompaniment energy $\langle a\rangle_i=\norm{a_i}^2/N$, so
\begin{equation}
\boxed{\ \ell_i(f)\approx\ell_{\text{clean}}+c\,\eps\,g(\langle a\rangle_i),\quad g\
\text{increasing}\ }
\end{equation}
per-chunk irreducible loss is monotone in accompaniment energy. Keeping the lowest-loss
chunks therefore keeps low-$\langle a\rangle$ chunks and drops high-$\langle a\rangle$ ones
--- the low-$\langle a\rangle$ chunks are the ones whose $\eps a$ bleed is quietest, the
\emph{effectively cleanest} targets. Trimming is \textbf{curriculum-by-cleanliness}: each
step trains on the least-contaminated slice of an entirely-contaminated dataset.

\paragraph{Testable signature.}
\begin{equation}
\boxed{\ \keptbar < \dropbar\ }
\end{equation}
the kept chunks' mean accompaniment energy is below the dropped chunks'. The loop logs this
every 500 steps (raw same-track $\langle a\rangle_i$ at the vocal window --- the $\eps$- and
augmentation-invariant covariate). No separation $\Rightarrow$ no cleanliness gradient
$\Rightarrow$ H-06b refuted with a mechanism.

\section{Statistics}

Pooled between-seed noise over the three 3-seed cells:
$\sigma_{\text{seed}}=\sqrt{\tfrac13\sum_c\hat\sigma_c^2}$. Recovery fraction
$\rho=(s_{\text{trim}}-s_{\text{bleed}})/(s_{\text{clean}}-s_{\text{bleed}})=u/d$; by the
delta method, treating the three 3-seed means as independent with variances
$\hat\sigma_c^2/n$,
\begin{equation}
\Var(\rho)\approx\Big(\tfrac1d\Big)^2\tfrac{\hat\sigma_{\text{trim}}^2}{n}
 +\Big(\tfrac{u}{d^2}\Big)^2\tfrac{\hat\sigma_{\text{clean}}^2}{n}
 +\Big(\tfrac{u-d}{d^2}\Big)^2\tfrac{\hat\sigma_{\text{bleed}}^2}{n},\quad
 \SE(\rho)=\sqrt{\Var(\rho)},
\end{equation}
CI $\rho\pm z\,\SE$. When $d\le0$ (H-06a unsupported), $\rho$ is not evaluable. Exactly two
pre-registered paired contrasts are tested on the 50 clean test tracks (bleed30$-$clean;
trim30$-$bleed30) with a bootstrap 95\% CI $+$ Wilcoxon; with only two primary pairs we
report both CIs rather than a family-wise correction.

\bigskip
\noindent\emph{Cross-reference.} \S1 \code{data/corrupt.py}; \S2 \code{losses/l1mag.py};
\S3 \code{analysis/bleed.py::prediction\_line} (machine-checked in
\code{tests/test\_bleed\_analysis.py}); \S4--5 \code{losses/trimmed.py} $+$ the
\code{acc\_energy} telemetry in \code{data/musdb\_dataset.py}$\to$\code{train/loop.py};
\S6 \code{analysis/bleed.py::recovery\_fraction} $+$
\code{analysis/scaling.py::pooled\_seed\_sigma}.

\end{document}
